\section{Choosing the unit and calendar weights together}
\label{exec17:s6-joint}

The block comparison above evaluates particular choices of weights. A bank
designing a comparison faces a further question: should it improve the donor
balance, the calendar balance, or both? The two choices interact. A residual
factor mismatch is a product of their imbalances, while independent-cell
noise is multiplied by both squared weight norms. Optimizing one axis while
holding the other fixed need not reveal the best joint choice. We can solve
one small reference problem completely, including its global alternatives,
without claiming that an estimated synthetic-control routine solves that
oracle problem.~\citep{exec17_s6_extension}

Use a new symmetric factor table, rather than change the inputs of the earlier
comparison. There is one treated region with loading 2, two donors with
loadings 0 and 1, one post period with factor 2, and two pre periods with
factors 0 and 1. The untreated mean in each cell is the product of its loading
and factor, plus unrestricted additive unit and time effects. The outcome
errors in the nine cells are independent, centered and have common known
variance $s>0$. A constant treatment effect $\tau$ enters only the treated
post cell. All choices below therefore estimate that same $\tau$; they do
not change the population or horizon to improve a score.

Put donor weights $(1-a,a)$ on the two donors and pre-period weights
$(1-b,b)$ on the two pre periods. Both pairs sum to one, even when a weight
is negative. In row order treated, donor zero, donor one, and column order
post, pre zero, pre one, the signed contrasts are
\begin{equation}
 q=(1,a-1,-a)',\qquad r=(1,b-1,-b)',\qquad
 \widehat\tau=\sum_{i=1}^3\sum_{j=1}^3q_i r_jY_{ij}.
 \label{exec17:s6-joint-contrast}
\end{equation}
Each contrast sums to zero, so the additive unit and time effects cancel.
The coefficient of the treated post cell is one. The remaining factor bias
is $(2-a)(2-b)$. Define $Q(x)=1+(1-x)^2+x^2$, the squared norm of one
axis's contrast. Independence of the nine cell errors gives the exact risk
\begin{equation}
 {\cal R}_s(a,b)
  =(2-a)^2(2-b)^2+sQ(a)Q(b).
 \label{exec17:s6-joint-risk}
\end{equation}
The first term is squared bias for $\tau$, and the second is its variance.
This expression is conditional on the declared factors and fixed weights.
It is a mean squared error, with squared outcome units. It is neither a
confidence radius nor a realized profit. Monetary scaling of an estimation
error would need the exposure and decision loss to be specified separately.

For any fixed $b$, this is a strictly convex quadratic in $a$. Its derivative
and unique conditional minimizer are
\begin{align}
 \frac{1}{2}\frac{\partial{\cal R}_s}{\partial a}
  &=(a-2)(2-b)^2+sQ(b)(2a-1),\nonumber\\
 a^*(b)&=\frac{2(2-b)^2+sQ(b)}{(2-b)^2+2sQ(b)}.
 \label{exec17:s6-joint-conditional}
\end{align}
The denominator is positive for $s>0$. There is a corresponding expression
for $b^*(a)$ by symmetry. When the other axis is already exact, for example
$b=2$, the optimum is $a=1/2$: factor bias is already zero, so the remaining
axis minimizes its noise norm. Matching the factor exactly on both axes
spends variance twice without reducing bias below zero. Conversely, when
the other axis retains imbalance, moving $a$ toward the extrapolating value
2 can have a benefit. Separate coordinatewise convexity does not establish
joint convexity or the uniqueness of a joint solution.

\subsection{A global certificate and two equally good comparisons}

Introduce $x=a-1$, $y=b-1$, $p=x+y$ and $t=xy$. The factor bias is
$1-p+t$. Direct multiplication gives
$Q(a)Q(b)=4(1+p+p^2-t+pt+t^2)$. Completing squares in these symmetric
coordinates yields
\begin{equation}
 \begin{split}
 {\cal R}_s(a,b)
  ={}&\frac{9s}{1+s}+3s(p+t)^2\\
     &+(1+s)\left\{p-t-\frac{1-2s}{1+s}\right\}^2.
 \end{split}
 \label{exec17:s6-joint-certificate}
\end{equation}
This identity is a lower-bound certificate for every real $a,b$. It does
not rely on a numerical search or on convergence of alternating updates.
For $0<s\leq1/2$, both squares can vanish with
\begin{equation}
 p=\frac{1-2s}{2(1+s)},\qquad t=-p,\qquad
 x,y=\frac{p\mathbin{\pm}\sqrt{p^2+4p}}{2}.
 \label{exec17:s6-joint-global}
\end{equation}
Here $p\geq0$, so the two roots are real. Substituting them produces actual
unit and calendar weights attaining the lower bound $9s/(1+s)$. Feasibility
of these roots matters: completing squares in $p,t$ alone would not prove
attainment by real $x,y$ for every noise level. No conclusion for $s>1/2$
is inferred from this zero-square construction.

At $s=1/8$, $p=1/3$ and $t=-1/3$. There are two distinct global solutions:
\begin{equation}
 (a,b)=\left(\frac{7+\sqrt{13}}6,\frac{7-\sqrt{13}}6\right)
 \quad\hbox{or the reversed pair},\qquad
 {\cal R}_{1/8}=1.
 \label{exec17:s6-joint-numeric}
\end{equation}
One coordinate is about 1.7676 and the other about 0.5657. The first
solution extrapolates across donor regions while using positive pre-period
weights; the second extrapolates across calendar periods while using
positive donor weights. Neither axis is exactly balanced. Equal risk comes
from the reference table's symmetry and common cell variance, rather than
from a general equivalence between borrowing information across regions
and borrowing it across time. A bank with common calendar shocks or unequal
exposures would usually break that symmetry.

At the boundary $s=1/2$, the certificate has only one real minimizer,
$a=b=1$, with risk 3. To see uniqueness directly, at $a=1+x,b=1+y$,
\begin{equation}
 {\cal R}_{1/2}(1+x,1+y)-3
   =3(x+y)^2+3x^2y^2.
 \label{exec17:s6-joint-boundary}
\end{equation}
Equality requires both $x+y=0$ and $xy=0$, hence $x=y=0$. The pair of
asymmetric optima at smaller positive noise therefore merges at this
boundary. This is a property of this finite reference objective. It is
not an asymptotic phase transition established for a fitted estimator.

\subsection{What constraints and exact matching cost}

Suppose policy requires nonnegative weights on both axes. Then
$(a,b)\in[0,1]^2$. At $s=1/8$, the first derivative in
\eqref{exec17:s6-joint-conditional} is strictly negative throughout the
square. Indeed $(a-2)(2-b)^2\leq-1$, while
$sQ(b)(2a-1)\leq1/4$, since $Q(b)\leq2$ on that interval. Thus
$\partial{\cal R}_{1/8}/\partial a\leq-3/2$. The same bound holds for
the other derivative. Increasing either coordinate until it reaches one
strictly lowers risk, proving the constrained global optimum $(1,1)$.
It is not merely the best point found on a grid.

The four comparisons on the same nine-cell law are
\begin{center}
\begin{tabular}{lrrr}
\toprule
Choice & $a$ & $b$ & Mean squared error\\
\midrule
Best nonnegative pair & $1$ & $1$ & $3/2$\\
Exact unit balance, minimum time norm & $2$ & $1/2$ & $9/8$\\
Exact balance on both axes & $2$ & $2$ & $9/2$\\
Unrestricted joint optimum & \multicolumn{2}{c}{pair in \eqref{exec17:s6-joint-numeric}} & $1$\\
\bottomrule
\end{tabular}
\end{center}
The exact-unit solution is also the best choice with $a$ constrained to 2:
its zero bias leaves $6sQ(b)$, minimized at $b=1/2$. Requiring both axes
to fit exactly multiplies two norms of 6 and gives $36s=9/2$. For the
nonnegative pair the residual bias is one, and the variance is $4s=1/2$.
The joint optimum accepts some residual bias to reduce the total risk below
the exact-unit alternative. These identities explain the comparisons;
observing a small pre-period residual alone would not do so.

There is a separate direct check. Each of the nine noise coefficients is
$q_i r_j$. Enumerating all $2^9=512$ independent signs for standardized
cell errors gives zero mean and second moment
$\sum_{i,j}q_i^2r_j^2$. Adding the squared factor-table contrast recovers
the displayed risks. For the irrational weights in
\eqref{exec17:s6-joint-numeric}, multiplying their exact cell coefficients
and summing squared coefficients also gives risk one. Thus neither a
floating-point optimizer nor a rounded square root determines the ranking.
\citep{exec17_s6_extension}

The result suggests a useful prototype diagnostic: retain both weight
vectors and the declared noise law when investigating a discrepancy between
two synthetic comparisons. It does not prescribe the oracle weights for an
actual bank. Its post factor and loadings are known here; a real untreated
post factor is counterfactual. Learning the weights from noisy pre outcomes
can correlate them with those errors. A contamination class, common shocks,
missing cells or a constraint on extreme weights changes the objective.
The earlier signed-contamination and calendar calculations show concrete
ways that the apparent gain can disappear. An estimated procedure needs its
own target-preserving selection rule and uncertainty argument. The global
certificate supplies an exact reference for a deliberately restricted law,
with which such a future implementation can be compared.

\section{Valuing replacement of the entire feedback rule}
\label{exec17:s6-full-rule}

The response to a new action innovation within a regime leaves a different
investment question unresolved. A bank may propose replacing the old
feedback rule for an already occupied portfolio, without adding an action
innovation. We now compare those entire regimes from the same nonzero
initial state, charging the actions each would actually take. An apparently
accurate model of the old path still leaves an interval of replacement
values, because the observed controller can hide the separate persistence
and action-effect parameters.

Retain the structural outcome law in \eqref{exec17:s6-feedback-law}, with
future structural disturbances and action innovations zero. Set $Y_{-1}=1$,
$k=2/5$ and the observed reduced coefficient $\lambda=9/10$. Supply the
continuous sensitivity restriction
\begin{equation}
 \beta\in[1,2],\qquad \rho=\lambda-k\beta.
 \label{exec17:s6-controller-class}
\end{equation}
Every world in this class gives exactly the same old state and action path,
for any length of observation under $A_t=kY_{t-1}$. The line
$\rho+k\beta=\lambda$ follows from the old deterministic observation law;
the numerical interval $[1,2]$ is additional supplied information. The
observations do not select that interval, supply probabilities over it or
make it a confidence region. Its use is a sensitivity analysis conditional
on a declared restriction and on structural invariance after replacement.

Replace $k$ by a fixed $k'$ at period zero. Define
$a(\beta)=\lambda+(k'-k)\beta$. Direct recursion gives
\begin{align}
 Y_t^{\rm new}&=a(\beta)^{t+1},&
 A_t^{\rm new}&=k'a(\beta)^t,\nonumber\\
 Y_t^{\rm old}&=\lambda^{t+1},&
 A_t^{\rm old}&=k\lambda^t,\qquad t=0,1,\ldots.
 \label{exec17:s6-full-paths}
\end{align}
Both first actions use the common initial state. There is no impact
innovation of size one and no factor $\beta$ multiplying the initial
state. Those features distinguish these paths from the within-regime
innovation responses developed earlier. If the initial state were zero
under these zero-disturbance premises, both paths would remain zero and
replacement would deliver only its setup loss.

Let $d\in(0,1)$ discount one calendar period, $m$ be dollars of margin per
state-period, $q>0$ be dollars per squared action-period, and $K$ be the
one-time replacement cost. The full paired-regime incremental value through
index $H$ is
\begin{equation}
 \begin{split}
 V_H(\beta;k')={}&\sum_{t=0}^H d^t
 \left[m\{a(\beta)^{t+1}-\lambda^{t+1}\}\right.\\
 &\left.\hspace{20mm}-q\{(k')^2a(\beta)^{2t}-k^2\lambda^{2t}\}\right]-K.
 \end{split}
 \label{exec17:s6-full-finite}
\end{equation}
The old action cost is subtracted from the new action cost, so a saving in
old operating expense raises incremental value. Omitting that old expense
would instead price the new policy against a costless comparator. The
summation contains $H+1$ periods; the setup cost is paid once. The units and
initial state have been normalized explicitly rather than extrapolated
from a local response to an unspecified portfolio.

For $G_H(z)=\sum_{t=0}^Hz^t$, equal to $(1-z^{H+1})/(1-z)$ when
$z\ne1$ and to $H+1$ when $z=1$, the finite formula becomes
\begin{equation}
 \begin{split}
 V_H={}&m\{aG_H(da)-\lambda G_H(d\lambda)\}\\
       &-q\{(k')^2G_H(da^2)-k^2G_H(d\lambda^2)\}-K.
 \end{split}
 \label{exec17:s6-full-geometric}
\end{equation}
This expression is finite for every finite $H$, including a rule whose
infinite discounted cost fails to converge. It makes a short explicit
planning horizon a substantive alternative to assuming a stationary
infinite future. The finite target still needs a defended terminal
treatment if the bank retains assets or liabilities beyond $H$.

\subsection{Robust stability and finite discounted cost are different}
\label{exec17:s6-rule-boundaries}

State stability requires $|a(\beta)|<1$. Since $a$ is affine in $\beta$,
its largest absolute value over $[1,2]$ is attained at an endpoint. Solving
both endpoint inequalities, with the old values $k=2/5$ and
$\lambda=9/10$, gives the robust stability interval
\begin{equation}
 -\frac{11}{20}<k'<\frac9{20}.
 \label{exec17:s6-robust-stability}
\end{equation}
For example, the upper endpoint comes from the $\beta=2$ world:
$9/10+2(k'-2/5)<1$. The corresponding $\beta=1$ bound is weaker.
The lower endpoint follows in the same way by bounding this coefficient
below by $-1$. Equality gives a nondecaying state in an admissible world,
so neither endpoint is included.

Convergence of discounted gross benefits only needs $|da|<1$.
For nonzero $k'$ and positive $q$, convergence of discounted quadratic
action cost instead needs $da^2<1$. The latter is stricter for
$d\in(0,1)$. Write $r_d=d^{-1/2}$. For this class, $r_d>\lambda>0$,
and the endpoint intersection for a uniformly finite discounted value is
\begin{equation}
 k-\frac{r_d+\lambda}{2}<k'<
 k+\frac{r_d-\lambda}{2}.
 \label{exec17:s6-robust-cost}
\end{equation}
The old policy already satisfies its convergence conditions. At
$d=19/20$, the new interval is approximately
$(-0.562989,0.462989)$. It is wider than the robust state-stability
interval. Discounting can make the sum finite while a state grows;
a finite valuation is consequently insufficient as an operational
admissibility criterion. Feasible capital, funding and action limits would
need to be imposed before calling such a rule implementable.

The cost condition cannot be replaced by convergence of the benefit
series. Consider $k'=47/100$. Then
\begin{equation}
 a(\beta)\in[97/100,26/25],\qquad
 d\max a=247/250<1,\qquad
 d(\max a)^2=3211/3125>1
 \label{exec17:s6-rule-divergence}
\end{equation}
at $d=19/20$. All admissible worlds have convergent discounted benefits,
but worlds with
$\beta\geq\{\sqrt{20/19}-9/10\}/(7/100)$ have a nonconvergent
positive action-cost series. In each of those worlds, $V_H$ tends to
$-\infty$ as $H$ grows: the benefit and old-cost totals have finite
limits, whereas the new quadratic expense does not. At equality the
discounted expense remains constant per period, which also makes its
sum diverge. A finite positive number obtained by formally substituting
into a geometric formula outside its convergence range would be invalid.

Where all four relevant series converge, the infinite value is instead
\begin{equation}
 \begin{split}
 V_\infty(\beta;k')={}&m\left\{\frac{a}{1-da}
                       -\frac{\lambda}{1-d\lambda}\right\}\\
 &-q\left\{\frac{(k')^2}{1-da^2}
                       -\frac{k^2}{1-d\lambda^2}\right\}-K.
 \end{split}
 \label{exec17:s6-full-infinite}
\end{equation}
The restriction must accompany the reported value. It is not supplied by
the fact that the old path is stable, since observationally equivalent
old models can cross the boundary differently after the same rule change.

\subsection{Locating the adverse world rather than assuming it}
\label{exec17:s6-rule-value}

Even inside the finite-value region, a quadratic cost can make net value
nonmonotone. Evaluating the two endpoint worlds alone therefore does not
generally give a robust value. For fixed $k'$, differentiating the
finite-valued infinite-horizon expression with respect to $a$ gives
\begin{equation}
 \frac{\partial V_\infty}{\partial a}
  =\frac{m}{(1-da)^2}
   -\frac{2dq(k')^2a}{(1-da^2)^2}.
 \label{exec17:s6-full-derivative}
\end{equation}
The old comparator and setup expense do not depend on $a$. Their absence
from this derivative does not justify omitting them from the level of
value. The negative term can dominate near a squared-cost boundary;
the worst case needs a sign argument or consideration of interior extrema.

For a concrete pair of stable rules, set $m=100$, $q=10$, $K=200$ and
$d=19/20$, and compare $k'=21/50$ with $k'=11/25$. Across both rules
and every $\beta\in[1,2]$, $23/25\leq a\leq49/50$. A single analytic
lower bound is
\begin{equation}
 \begin{split}
 \frac{\partial V_\infty}{\partial a}
 &\geq\frac{100}{\{1-(19/20)(23/25)\}^2}\\
 &\quad-\frac{2(19/20)10(11/25)^2(49/50)}
                   {\{1-(19/20)(49/50)^2\}^2}
 >5829.
 \end{split}
 \label{exec17:s6-controller-monotone}
\end{equation}
The first term uses its smallest value over the interval. The magnitude
of the second term uses its largest value: for $a\geq0$ below the
cost boundary, $a/(1-da^2)^2$ is increasing. Both rules have $k'>k$,
so $a$ also increases with $\beta$. These facts prove that the smallest
infinite value is at $\beta=1$ for each rule, rather than assume that
the largest action effect is always the adverse world.

The same endpoint is adverse at every finite $H$ in this example. For
$t\geq1$, the derivative of the new period's discounted contribution is
\begin{equation}
 d^t a^t\{m(t+1)-2tq(k')^2a^{t-1}\}>0.
 \label{exec17:s6-finite-monotone}
\end{equation}
Here $0<a\leq1$ and $m>2q(k')^2$ suffice; at $t=0$ the derivative
is $m>0$. Thus each summand, and their finite sum, is increasing in
$a$. This finite argument does not depend on an exchange of a derivative
and an infinite series. With a larger cost, a different action range or
a state above one, its sufficient bound would need to be reassessed.

The exact rational values, rounded to cents only for display, are
\begin{center}
\begin{tabular}{rrrrrr}
\toprule
$k'$ & $\beta$ & $a$ & $V_{20}$ & $V_{80}$ & $V_\infty$\\
\midrule
$0.42$ & $1$ & $0.92$ & $-112.57$ & $-92.60$ & $-92.59$\\
$0.42$ & $2$ & $0.94$ & $-4.44$ & $53.68$ & $53.77$\\
$0.44$ & $1$ & $0.94$ & $-5.48$ & $52.61$ & $52.70$\\
$0.44$ & $2$ & $0.98$ & $294.31$ & $580.12$ & $584.45$\\
\bottomrule
\end{tabular}
\end{center}
Values are dollars per normalized initial exposure, including both regimes'
operating expenses and the USD 200 setup cost. Under the supplied class,
rule 0.42 cannot be declared profitable in every world even at an infinite
horizon. Rule 0.44 has positive worst-case infinite value, but its
worst-case value through index 20 remains negative. The table's endpoints
are rigorous extrema for these rules because of the preceding derivative
bounds. They are not percentiles or probabilities of gain.

Notice that the two middle rows share $a=0.94$, but their values differ.
The rule coefficients are different, so their action expenses differ even
when their new state paths coincide. This is a concrete reason to retain
the action path when translating a predicted state improvement into money.
A model returning only a state forecast could hide the operating-cost
distinction. Similarly, a comparison of $V_{20}$ with $V_\infty$ changes
the investment target; it is not a numerical refinement of the same
finite-horizon quantity.

\subsection{A tail certificate for the whole paired comparison}

The earlier innovation tail cannot be reused unchanged for this question.
Subtracting \eqref{exec17:s6-full-geometric} from
\eqref{exec17:s6-full-infinite} leaves four tails:
\begin{equation}
 \begin{split}
 V_\infty-V_H={}&m\left\{
 \frac{a(da)^{H+1}}{1-da}
 -\frac{\lambda(d\lambda)^{H+1}}{1-d\lambda}\right\}\\
 &-q\left\{\frac{(k')^2(da^2)^{H+1}}{1-da^2}
 -\frac{k^2(d\lambda^2)^{H+1}}{1-d\lambda^2}\right\}.
 \end{split}
 \label{exec17:s6-full-tail}
\end{equation}
For the positive coefficients in the two-rule example, the triangle
inequality gives an absolute bound by adding these four positive
magnitudes. Maximizing the new-state and action terms at $a=0.98$ and
$k'=0.44$ yields a uniform bound of USD 4.3532, below USD 4.36, at $H=80$.
The actual largest endpoint tail is about USD 4.3231; cancellation makes
it smaller than the four-term bound. A uniformly safe truncation allowance
should use the defended envelope, rather than an observed cancellation
at a convenient world.~\citep{exec17_s6_extension}

Directly evolving old and new states, taking their separate actions and
accumulating discounted gains reproduces the finite geometric formula.
Adding the exact remaining tails reproduces the infinite value, for both
endpoint worlds and an interior world with $\beta=3/2$. These checks also
retain the first action and its cost. They would fail if the calculation
were silently replaced by a within-regime innovation response or if the
initial state were set to zero.

The decision implication remains conditional. A positive worst-case value
for rule 0.44 is warranted only for this supplied structural class, action
feasibility, initial exposure, cost law and horizon. The old observations
cannot certify the class or structural invariance. An exogenous action
perturbation, a credible restriction on persistence, or a justified
controlled pilot could narrow the class; a public panel can supply
realistic clocks and exposures without supplying that missing causal
truth. These are concrete next information choices. Until they are
resolved, the proposal should return the class of values and the
convergence restrictions alongside any chosen policy, so that an
attractive forecast is not mistaken for an identified bank investment.
